\section{¿Puede existir una transferencia positiva entre la comprensión y la generación?}

Se han presentado anteriormente muchos «modelos unificados», pero la unificación de la arquitectura no implica transferencia de capacidades. Un modelo puede utilizar el mismo backbone para realizar VQA y generación de imágenes, pero aún así requerir dos conjuntos de representaciones; también puede generar videos realistas sin poseer un mejor razonamiento espacial. En la segunda mitad de la clase se puso esta incertidumbre en el centro y se propuso que las diferencias en la estructura de información de las señales lingüísticas y visuales podrían ser una causa raíz.

\subsection{Primero, clarifiquemos la dirección de la transferencia}

Todos los «modelos unificados» presentados anteriormente pueden llevar a suponer que las tareas se ayudan mutuamente. En esta sección corregimos esa intuición: la transferencia tiene dirección y también requiere una línea base de tarea única emparejada; por lo tanto, las diapositivas 047--048 se encargan de definir los experimentos en lugar de anunciar respuestas.

\lecturefigure{slide-047.jpg}{Segunda línea de investigación: si la comprensión y la generación tienen una transferencia positiva}{Diapositiva oficial 47; Clase 00:35:11--00:35:30}

\begin{importantbox}{Definición verificable de la transferencia positiva}
Tras controlar el volumen de datos, la capacidad de cómputo, los parámetros y el protocolo de evaluación, solo podemos denominar a la tarea $A\rightarrow B$ como una transferencia positiva si el entrenamiento conjunto hace que el rendimiento de la tarea objetivo $B$ sea superior al de la línea base que entrena únicamente $B$. Si el aumento es simplemente debido a un modelo total más grande, un entrenamiento más prolongado o el uso de anotaciones adicionales, no se puede atribuir la ganancia a la transferencia.
\end{importantbox}

\lecturefigure{slide-048.jpg}{Dos caminos: Comprensión $\rightarrow$ Generación y Generación $\rightarrow$ Comprensión}{Diapositiva oficial 48; Clase 00:35:30--00:36:10}

\begin{knowledgebox}{Lectura de la figura: las dos flechas requieren experimentos diferentes}
La hipótesis del lado izquierdo (de la comprensión a la generación) asume que la percepción, el etiquetado o la representación del razonamiento pueden mejorar la fidelidad del prompt, el diseño de layout y la edición. La hipótesis del lado derecho (de la generación a la comprensión) asume que el modelado del mundo generativo ayuda al modelo a dominar objetos, movimiento o causalidad. Para verificar la flecha izquierda, se debe fijar el generador e incorporar entrenamiento de comprensión; para verificar la flecha derecha, hay que fijar el benchmark de comprensión e incorporar objetivos de generación. No se pueden usar un mismo conjunto de muestras cualitativas para demostrar ambas direcciones simultáneamente.
\end{knowledgebox}

\begin{warningbox}{Compartir parámetros no es prueba de transferencia}
Que dos tareas pasen por el mismo Transformer solo indica que podrían interactuar. Si hay conflicto de gradientes, desequilibrio en la escala de pérdida o incompatibilidad de representaciones, el entrenamiento conjunto podría incluso provocar una transferencia negativa. Los resultados de MoT ya proporcionan un límite contraejemplo: hubo ganancias en la especialización para la generación de imágenes, pero no se observó un aumento correspondiente en la comprensión de imágenes.
\end{warningbox}

\subsection{¿Por qué el lenguaje parece más una «interfaz cognitiva comprimida»?}

Los tokens del lenguaje a menudo ya han sido seleccionados por humanos: una frase conserva objetos, relaciones, intenciones y anomalías, dejando fuera gran parte de los detalles pixelales. Las imágenes y los videos se acercan más a la observación sensorial cruda e incluyen iluminación, textura, fondo, movimiento de cámara y fotogramas repetidos. Por lo tanto, la misma predicción del siguiente paso tiene diferentes ratios señal-ruido en ambos tipos de datos.

\lecturefigure{slide-051.jpg}{Compresión lingüística, observación visual pasiva, geometría de pérdida y redundancia de fotogramas}{Diapositiva oficial 51; Clase 00:36:10--00:39:11}

\begin{knowledgebox}{Lectura de la figura: los cuatro burbujas forman una hipótesis causal, no una ley demostrada}
El lado izquierdo señala que el lenguaje es una abstracción cognitiva altamente comprimida; el lado derecho lista tres dificultades de las imágenes/videos: observación pasiva, paisaje de pérdida y correlación de fotogramas. Juntas explican por qué la pérdida visual puede ser baja pero la generación aún no natural, y también explican por qué los videos poseen enormes bytes pero no necesariamente proporcionan supervisión independiente equivalente. La figura apoya las hipótesis de investigación, pero aún no demuestra que el lenguaje sea naturalmente superior a todas las representaciones visuales.
\end{knowledgebox}

Si la correlación entre fotogramas adyacentes del video es aproximadamente $\rho$, una intuición común de AR(1) es:
$$
N_{\text{eff}}\approx N\frac{1-\rho}{1+\rho}.
$$
$N$ es el número de fotogramas y $N_{\text{eff}}$ es el número aproximado de muestras independientes. Cuanto más cerca esté $\rho$ de 1, menos información independiente aportará cada nuevo fotograma. Esta fórmula es solo una intuición estadística y no debe usarse directamente para videos complejos; nos recuerda que «más fotogramas» no equivale a «más eventos semánticos».

\begin{warningbox}{Una baja pérdida perceptual no garantiza alta calidad humana}
El MSE de píxeles/latentes promedia múltiples futuros posibles, y la difusión optimiza en el espacio del ruido. Algunos errores son extremadamente significativos para la visión humana pero representan un valor numérico muy pequeño; otros cambios de píxeles tienen un valor numérico grande pero no cambian el significado. Por lo tanto, es necesario combinar preferencias humanas, relaciones de objetos, consistencia temporal y éxito de tareas, en lugar de perseguir únicamente la pérdida de entrenamiento más baja.
\end{warningbox}

\teachervoice{Del 00:38:11 al 00:38:21, el docente señaló que incluso si la generación de imágenes/videos aún no se ve bien para los humanos, la pérdida puede ya estar mejorando. Esta advertencia separa claramente la métrica de optimización de la calidad visible por el usuario.}

\subsection{El procesamiento multimodal digital aún no es inteligencia física}

La sección anterior explicó las dificultades de las señales de aprendizaje visual; esta sección amplía el alcance desde la optimización hasta la capacidad del sistema. La diapositiva 052 busca responder: aunque un modelo pueda procesar y generar diversos contenidos digitales, ¿qué capas faltan para lograr una percepción continua, control y acción?

\lecturefigure{slide-052.jpg}{De procesamiento de información digital a inteligencia multimodal más amplia en el mundo físico}{Diapositiva oficial 52; Clase 00:39:12--00:40:04}

\begin{knowledgebox}{Lectura de la figura: el modelo central solo cubre una parte de la imagen completa}
El lado izquierdo muestra información digital (texto, imagen, video, audio); el centro lista Chameleon, Transfusion, MoT y generación cruzada modal; el lado derecho del mundo físico también requiere conciencia espacial, comprensión del mundo en tiempo real, control, exploración e interacción a largo plazo. La conclusión inferior es contenida: los modelos actuales sobresalen en procesamiento de información multimodal, pero están muy lejos de la inteligencia física completa del mundo real.
\end{knowledgebox}

El ciclo físico se puede escribir como:
$$
o_t\xrightarrow{\text{encode}}h_t
\xrightarrow{\text{reason/plan}}a_t
\xrightarrow{\text{environment}}o_{t+1}.
$$
La generación digital suele evaluar solo $a_t$ (el contenido de salida); los sistemas robóticos también deben evaluar si la transferencia del entorno es segura, si la estimación de estado se está desviando, si el retardo cumple con el ciclo de control y si hay recuperación tras un fallo.

\begin{importantbox}{¿Qué falta para pasar de un modelo omni a una inteligencia encarnada?}
No falta solo más modalidades, sino percepción activa, estimación de estado, memoria espaciotemporal, control en bucle cerrado, calibración de incertidumbre, infraestructura en tiempo real y restricciones de seguridad. Enviar una imagen a un LLM y que un robot vea, escuche y actúe continuamente difieren en sus requisitos de sistema por un ciclo completo cerrado.
\end{importantbox}

\subsection{Resumen de la sección}

La transferencia entre comprensión y generación debe verificarse por dirección, línea base y cómputo emparejado. El lenguaje proporciona señales abstractas de alta densidad; la visión proporciona datos sensoriales de alto ancho de banda pero con alta redundancia; sus mejores representaciones y pérdidas no necesariamente son las mismas. Los modelos omni actuales han progresado significativamente en el procesamiento de información digital, pero aún no cubren el ciclo continuo de estado y acción requerido por el mundo físico.


